\paragraph{Calidad}
\textbf{Hallazgo.} Montgomery entrega siempre 4020x4892 pixeles con independencia de cuanto colimo el tecnico. El area sin senal llega al 67 por ciento en el caso mas colimado, frente a una mediana del 1 por ciento en Shenzhen.

\textbf{Evidencia.} fraccion de pixeles por debajo de 8: mediana 19.4 por ciento en Montgomery y 1.1 por ciento en Shenzhen sobre 562 imagenes de entrenamiento; 31 imagenes superan el 25 por ciento, todas de Montgomery

\textbf{Decision.} Delimitar el campo irradiado como region binaria, sin forzarlo a rectangulo y sin rellenar huecos, y usarlo como region de interes para histogramas, umbralizacion y denominador de la fraccion de area

\textbf{Efecto esperado.} Las medidas dejan de depender de la apertura del colimador. El fondo dentro de la region pasa del 20-32 por ciento que dejaba el recorte rectangular al 0 por ciento

\textbf{Limitacion.} Las axilas de un paciente delgado se registran al nivel del fondo y quedan fuera de la region. Es coherente con medir el cuerpo del paciente, pero significa que la region no es el torax completo

\paragraph{Escala y roi}
\textbf{Hallazgo.} Las mascaras de referencia de Shenzhen estan a 512x512 mientras que sus imagenes rondan 2940x3000, y contienen valores intermedios propios del antialiasing del borde

\textbf{Evidencia.} 388 de 486 mascaras requieren reescalado; los valores intermedios son el 0.14 por ciento de los pixeles y binarizar en cualquier umbral entre 1 y 255 cambia el area menos del 0.3 por ciento

\textbf{Decision.} Reescalar la mascara a la imagen por vecino mas cercano y binarizar en 128. Reducir despues imagen y mascara al mismo lado de trabajo de 512 pixeles antes de segmentar

\textbf{Efecto esperado.} El Dice se calcula sobre la misma rejilla para ambos origenes, y el tiempo de proceso baja de horas a 26 minutos para las 486 imagenes

\textbf{Limitacion.} El reescalado no crea detalle que la anotacion no tenga: la mascara de Shenzhen conserva la precision de 512x512 aunque se represente sobre una rejilla mayor

\paragraph{Caracteristicas}
\textbf{Hallazgo.} La simetria izquierda-derecha y la compacidad son los descriptores mas inestables frente a la estrategia de segmentacion

\textbf{Evidencia.} error relativo mediano de la simetria entre 0.647 y 0.839 segun la estrategia, y de la compacidad entre 0.410 y 0.593; frente a 0.066 a 0.158 de la relacion de aspecto

\textbf{Decision.} Conservar los cinco descriptores y reportar su error relativo por separado, en lugar de promediarlos o descartar los inestables

\textbf{Efecto esperado.} La respuesta a la pregunta distingue que descriptores son reproducibles y cuales no, que es lo que el objetivo cuatro pide

\textbf{Limitacion.} La inestabilidad de la simetria proviene en parte de que una segmentacion que pierde un campo pulmonar produce asimetria maxima, de modo que el descriptor mezcla anatomia con fallo de metodo

\paragraph{Preprocesamiento}
\textbf{Hallazgo.} Ninguna estrategia de realce mejora la linea base sin realce de forma consistente en los dos origenes

\textbf{Evidencia.} de las seis combinaciones con realce, la mejor alcanza 4 de 5 descriptores mejorados en Shenzhen pero solo 2 de 5 en Montgomery; ninguna llega a 3 de 5 en ambos. El Dice mediano mas alto es de base\_adap con 0.810, sin realce

\textbf{Decision.} Seleccionar base\_adap como combinacion del flujo final y reportar el resultado negativo del realce como hallazgo, no como fallo

\textbf{Efecto esperado.} El flujo final es mas simple y mas reproducible entre sitios: base\_adap da 0.8077 en Montgomery y 0.8102 en Shenzhen

\textbf{Limitacion.} El resultado es valido para los parametros declarados de cada realce. Un CLAHE con otro limite de recorte o rejilla podria comportarse de otro modo; no se exploro ese espacio

\paragraph{Evaluacion futura}
\textbf{Hallazgo.} El efecto del sitio de adquisicion supera al de la clase: en intensidad los dos origenes forman nubes casi disjuntas, mientras que dentro de cada origen las clases apenas difieren

\textbf{Evidencia.} mediana de intensidad 111.0 en Montgomery y 173.5 en Shenzhen, con IQR de 151.5 frente a 90.0; dentro de Montgomery la mediana es 110.0 para TB negativo y 111.0 para TB positivo

\textbf{Decision.} Reportar toda metrica estratificada por origen y exigir consistencia en ambos como criterio de exito, nunca un promedio agregado

\textbf{Efecto esperado.} Una estrategia que funcione en un solo sitio no puede presentarse como solucion. El criterio ya descarto a clahe\_fijo, que agregado habria parecido una mejora

\textbf{Limitacion.} Con dos sitios no se puede distinguir si la diferencia proviene del detector, del protocolo o de la poblacion. Un tercer origen seria necesario para separarlos

\paragraph{Unidad de analisis}
\textbf{Hallazgo.} Las lecturas clinicas declaran dos grupos de imagenes del mismo paciente en Montgomery, y 90 imagenes comparten sexo y edad con alguna otra sin que eso pruebe identidad

\textbf{Evidencia.} dos referencias cruzadas explicitas en el texto de las lecturas; 34 combinaciones de sexo y edad con mas de una imagen, que afectan a 90 de las 138 imagenes de Montgomery

\textbf{Decision.} Agrupar solo los dos casos demostrados y declarar el resto como riesgo residual, en lugar de agrupar por sexo y edad

\textbf{Efecto esperado.} No se inventa independencia donde no se puede demostrar, ni se confunden pacientes distintos reduciendo artificialmente el numero de unidades

\textbf{Limitacion.} No se puede descartar que existan mas imagenes del mismo paciente sin declarar. La agrupacion por sexo y edad tampoco habria servido: uno de los dos pares demostrados tiene edades distintas porque el paciente cumplio anos entre estudios

